\section{GGPKD: Global Geometry through Local Alignment}
\label{sec:method}

GGPKD distills an embedding teacher into a smaller student by matching the
teacher's one-hop $k$NN transition rows. The method has only two
method-specific hyperparameters: the graph width \texttt{graph\_k} and the
auxiliary-row weight \texttt{row\_weight}. All other temperatures, weights,
candidate sets, and supervised rows are derived from the teacher graph. The
graph and teacher embeddings are cached before training and are not needed at
inference.

\subsection{Teacher graph}
\label{sec:graph}

Let $\mathcal D=\{x_i\}_{i=1}^{N}$ be the deduplicated training corpus. The
teacher encodes every text once, and its embeddings are $\ell_2$-normalized so
that all similarities are cosines. For each node $i$, we retrieve the
$k=\texttt{graph\_k}$ most similar non-self nodes. If $s_i^{(q)}$ is the
$q$-th largest retrieved cosine, the row bandwidth is
\begin{equation}
\tau_i=\frac{s_i^{(1)}-s_i^{(k)}}{\log k}.
\label{eq:row-bandwidth}
\end{equation}
Thus the $k$-th retrieved neighbor is $k$ times less likely than the nearest
neighbor before mutual filtering. Because the bandwidth scales with any affine
rescaling of the similarities, the resulting softmax row is affine-invariant
and requires no teacher-specific temperature tuning.

We retain an edge $i\rightarrow j$ only when $i$ and $j$ retrieve each other.
If mutual filtering isolates a node, its raw top-$k$ row is used as a fallback.
Writing the resulting neighbor set as $\mathcal N_i$, the teacher transition
row is
\begin{equation}
P_i^T(j)
=\frac{\mathbf 1[j\in\mathcal N_i]\exp(c_{ij}^T/\tau_i)}
{\sum_{u\in\mathcal N_i}\exp(c_{iu}^T/\tau_i)}.
\label{eq:teacher-transition}
\end{equation}
The bandwidth in Equation~\eqref{eq:row-bandwidth} is computed from the raw
top-$k$ scores before this filter, so it is defined from the same number of
neighbors for every node.

For numerical efficiency, each row retains the smallest highest-mass prefix
$\Omega_i\subseteq\mathcal N_i$ satisfying
\begin{equation}
\sum_{j\in\Omega_i}P_i^T(j)\geq 0.99,
\qquad
\widetilde P_i^T(j)
=\frac{P_i^T(j)}{\sum_{u\in\Omega_i}P_i^T(u)}.
\label{eq:row-truncation}
\end{equation}
The one-percent tolerance is a numerical-fidelity constant rather than a
method hyperparameter. It produces ragged rows while bounding the discarded
teacher mass.

\subsection{Deterministic candidate pool}
\label{sec:candidates}

The candidate set for anchor $i$ is exactly its truncated transition support
$\Omega_i$. GGPKD draws no hard or random negatives and performs no per-epoch
sampling: the candidates are a deterministic function of the graph and remain
fixed throughout training. Within a mini-batch $B$, candidates are deduplicated
into the shared pool
\begin{equation}
\mathcal P_B=\bigcup_{i\in B}\Omega_i,
\label{eq:shared-pool}
\end{equation}
so each distinct candidate text is encoded once. Ragged rows are padded with
their anchor index during collation; a self-mask removes that padding from all
softmaxes.

\subsection{Training objective}
\label{sec:objective}

The graph-scale loss matches each truncated teacher row with a student softmax
over the same columns and at the same stored bandwidth:
\begin{equation}
\mathcal L_{r=1}
=\frac{1}{|B|}\sum_{i\in B}
D_{\mathrm{KL}}\!\left(
\widetilde P_i^T
\,\middle\|\,
\operatorname{softmax}_{j\in\Omega_i}(c_{ij}^S/\tau_i)
\right).
\label{eq:graph-loss}
\end{equation}

This local row ranks candidates within an anchor's own neighborhood. To also
calibrate similarities across the batch, the ambient loss uses the complete
shared pool. Its temperature is derived from the graph as
$\tau_0=\operatorname{median}_i\tau_i$ and is shared by teacher and student:
\begin{equation}
\mathcal L_{r=0}
=\frac{1}{|B|}\sum_{i\in B}
D_{\mathrm{KL}}\!\left(
\operatorname{softmax}_{j\in\mathcal P_B\setminus\{i\}}(c_{ij}^T/\tau_0)
\,\middle\|\,
\operatorname{softmax}_{j\in\mathcal P_B\setminus\{i\}}(c_{ij}^S/\tau_0)
\right).
\label{eq:ambient-loss}
\end{equation}
The ambient and graph groups receive equal total weight; this balance is fixed,
not tuned.

GGPKD also reuses non-anchor candidates as auxiliary row centers. Let
$\mathcal J_B=\mathcal P_B\setminus B$ and expose the portion of row $j$ that is
already present in the shared pool:
\begin{equation}
\Omega_j^B=\mathcal N_j\cap\mathcal P_B,
\qquad
\mathcal J_B^+=\{j\in\mathcal J_B:|\Omega_j^B|\geq2\}.
\end{equation}
Batch anchors are excluded because $\mathcal L_{r=1}$ already supervises their
transition rows. The remaining rows are matched uniformly at their stored
bandwidths:
\begin{equation}
\mathcal L_{\mathrm{row}}
=\frac{1}{|\mathcal J_B^+|}\sum_{j\in\mathcal J_B^+}
D_{\mathrm{KL}}\!\left(
P_j^T\big|_{\Omega_j^B}
\,\middle\|\,
\operatorname{softmax}_{u\in\Omega_j^B}(c_{ju}^S/\tau_j)
\right),
\label{eq:row-loss}
\end{equation}
where the restricted teacher row is renormalized on $\Omega_j^B$. This term
requires no additional student encoding because all centers and columns are
already in $\mathcal P_B$.

The complete objective is
\begin{equation}
\mathcal L_{\mathrm{GGPKD}}
=\underbrace{\mathcal L_{r=0}+\mathcal L_{r=1}}_{\mathcal L_{\mathrm{rel}}}
+\lambda_{\mathrm{row}}\mathcal L_{\mathrm{row}}.
\label{eq:full-objective}
\end{equation}
The defaults are $\texttt{graph\_k}=200$ and
$\lambda_{\mathrm{row}}=\texttt{row\_weight}=1$. Increasing
\texttt{graph\_k} changes both neighborhood width and the bandwidth in
Equation~\eqref{eq:row-bandwidth}; those effects cannot be separated by a
single-axis sweep.

\paragraph{Cost.}
Exact graph construction requires $O(N^2d)$ similarity work, although an
approximate nearest-neighbor backend can replace the all-pairs search. During
training, the dominant cost is encoding the deduplicated shared pool;
$\mathcal L_{\mathrm{row}}$ reuses those embeddings. At deployment, inference
uses only the student encoder.
